\documentclass[11pt]{article}
\usepackage[margin=1.1in]{geometry}
\usepackage{amsmath,amssymb,amsthm}
\usepackage[T1]{fontenc}
\usepackage{lmodern}
\usepackage{hyperref}
\hypersetup{colorlinks=true,linkcolor=blue,urlcolor=blue,citecolor=blue}
\newtheorem{theorem}{Theorem}
\newtheorem{lemma}[theorem]{Lemma}
\newtheorem{proposition}[theorem]{Proposition}
\theoremstyle{definition}
\newtheorem{definition}[theorem]{Definition}
\setlength{\parskip}{2pt}
\title{A proof of the conjectured closed form for OEIS A278721}
\author{Adrian Perez Fontelles and Gaspard Moulinier, Independent researchers}
\date{09 September 2026}
\begin{document}
\maketitle
\section{The conjecture}

The Online Encyclopedia of Integer Sequences records \href{https://oeis.org/A278721}{A278721} as

\begin{quote}\itshape Number of nX2 0..3 arrays with rows in nondecreasing lexicographic order and columns in nonincreasing lexicographic order\end{quote}

and carries in its Formula section the line:

\begin{quote}\ttfamily Empirical: a(n) = (1/479001600)*n\textasciicircum{}12 + (1/7257600)*n\textasciicircum{}11 + (1/248832)*n\textasciicircum{}10 + (103/1451520)*n\textasciicircum{}9 + (12631/14515200)*n\textasciicircum{}8 + (2663/345600)*n\textasciicircum{}7 + (433049/8709120)*n\textasciicircum{}6 + (342073/1451520)*n\textasciicircum{}5 + (1269197/1555200)*n\textasciicircum{}4 + (3605419/1814400)*n\textasciicircum{}3 + (521131/166320)*n\textasciicircum{}2 + (997/360)*n + 1\end{quote}

That line carries no separate signature; the entry itself is by R. H. Hardin (Nov 27 2016).

The word {\itshape Empirical} --- equivalently {\itshape Conjecture} --- is the entry's own: the statement is recorded as fitted to the published terms, and no proof is given there. As of revision 5, last modified Nov 27 2016, the entry records no proof and links to none, so the statement is open. This note settles it.

Written out, the claim is

\[a(n) = 1 + \frac{997}{360}n + \frac{521131}{166320}n^{2} + \frac{3605419}{1814400}n^{3} + \frac{1269197}{1555200}n^{4} + \frac{342073}{1451520}n^{5} + \frac{433049}{8709120}n^{6} + \frac{2663}{345600}n^{7} + \frac{12631}{14515200}n^{8} + \frac{103}{1451520}n^{9} + \frac{1}{248832}n^{10} + \frac{1}{7257600}n^{11} + \frac{1}{479001600}n^{12}.\]

\section{The object}

The name is read as follows. The reading is not a guess: it is pinned against the entry's own published terms, and against an independent enumeration, before anything is proved (Section 7).

Let $A$ be an $n' \times 2$ array over $\{0,\dots,3\}$ with $n' = n$ rows.

Compare adjacent rows, and adjacent columns, as words in the lexicographic order. A {\itshape mistake} is an adjacent pair of rows that is not in nondecreasing order, or an adjacent pair of columns that is not in nonincreasing order; equal neighbours are never mistakes.

The condition is that there are exactly $0$ mistakes in all.

$a(n)$ is the number of such arrays. The entry's offset is 1, so its term of index $n$ is the count for $n$ rows.

\section{The count is a walk count}

Comparing two adjacent rows needs only those two rows. Comparing two adjacent columns needs the whole array, but only through one of three states per adjacent pair --- equal so far, already smaller, already larger --- and that state is updated one row at a time and never returns to `equal so far' once it has left it.

So take as state the last row, the $1$ column-comparison states, and the number of row mistakes so far capped at $1$. A state is accepting when the row mistakes plus the column mistakes its comparison states record come to exactly $0$.

An array of $n'$ rows is then exactly a walk of length $n'$ from the empty state to an accepting one, so $a(n) = w^{\mathsf T}M^{\,n'}f$ and the count is C-finite.

\section{The degree bound, derived}

The reachable part of the digraph has $68$ states, $17$ after trimming; the forward identification of states with equal numbers of completions of every length, then the mirror identification on the edges coming in, leaves $S = 13$ blocks, and the count satisfies the linear recurrence given by the characteristic polynomial of that $13\times13$ matrix.

\section{The decision procedure}

Let $c_1,\dots,c_D$ be the coefficients of a claimed recurrence $a(n) = \sum_{i=1}^{D} c_i\,a(n-i)$, let $q(t) = t^{D} - \sum_{i=1}^{D} c_i t^{D-i}$ be its characteristic polynomial, and put

\[u_m \;=\; \tilde w^{\mathsf T} L^{\,m-1} q(L)\,\tilde f \qquad (m \ge 1).\]

Expanding the powers of $L$ and using the displayed formula for $a$, one gets $u_m = a(n) - \sum_{i=1}^{D} c_i\,a(n-i)$ with $n = m + D$. So $u_m$ is exactly the residual of the claim at that index, and the recurrence holds at an index $n$ if and only if the corresponding $u_m$ vanishes.

Now $u_m = \tilde w^{\mathsf T} L^{m-1} y$ with $y = q(L)\tilde f$ a fixed vector, so $(u_m)$ satisfies the characteristic polynomial of $L$, of degree $13$: writing that polynomial as $t^{13} - \sum_{i=1}^{13} \lambda_i t^{13-i}$ gives $u_{m+13} = \sum_{i=1}^{13} \lambda_i u_{m+13-i}$. Hence $13$ consecutive vanishing residuals force every later residual to vanish, by induction. The test is therefore finite; and because it locates the {\itshape last} nonvanishing residual, it pins the threshold exactly rather than merely bounding it.

Everything is carried out in exact integer arithmetic on the vector $L^{m-1}y$. The matrix $M$ is never formed.

\section{Results}

\begin{theorem} For A278721, $a(n)$ is the polynomial $1 + \frac{997}{360}n + \frac{521131}{166320}n^{2} + \frac{3605419}{1814400}n^{3} + \frac{1269197}{1555200}n^{4} + \frac{342073}{1451520}n^{5} + \frac{433049}{8709120}n^{6} + \frac{2663}{345600}n^{7} + \frac{12631}{14515200}n^{8} + \frac{103}{1451520}n^{9} + \frac{1}{248832}n^{10} + \frac{1}{7257600}n^{11} + \frac{1}{479001600}n^{12}$ for every $n \ge 1$.\end{theorem}

{\itshape Proof.} A polynomial of degree $12$ is annihilated by the difference operator $(E-1)^{13}$, whose characteristic polynomial is $(t-1)^{13}$. So $a(n) - P(n)$, with $P$ the claimed polynomial, satisfies the linear recurrence whose characteristic polynomial is the product of that with the characteristic polynomial of $L$, of degree $S + 13 = 26$. Its values were computed exactly; all of them vanish, and more than $26$ consecutive vanishing values follow, so every later one vanishes as well. $\square$

\section{Verification}

Four checks were run, each reproducible from the code accompanying this note, and the first two are what pin the reading of the entry's English.

\begin{itemize}

\item {\bfseries The published terms.} The walk count reproduces all 26 terms the entry publishes, exactly, in integer arithmetic, with the index taken from the entry's stated offset (1) rather than fitted. The first of them are 10, 60, 275, 1050, 3492, 10401, 28288, 71266,~\dots

\item {\bfseries An independent enumeration.} For $n = 1,\dots,4$ every one of the $4^{2n}$ arrays was written out and tested cell by cell against the condition as the entry states it --- no transfer matrix, no row window, a separate piece of code. The counts agree with the walk count and with the entry.

\item {\bfseries The claim on the raw data.} Each claim was evaluated on the entry's published terms alone, with no model involved, wherever the proved range and the published data overlap. It holds there.

\item {\bfseries The annihilation test.} Carried to the derived bound $S = 13$ over the whole vector, in exact integer arithmetic, never sampled and never truncated.

\end{itemize}

\section{References}

\begin{enumerate}

\item OEIS Foundation Inc., \emph{The On-Line Encyclopedia of Integer Sequences}, entry \href{https://oeis.org/A278721}{A278721}, revision 5, last modified Nov 27 2016.

\item R. P. Stanley, \emph{Enumerative Combinatorics}, Vol.~1, 2nd ed., Cambridge University Press, 2012; \S4.7, the transfer-matrix method.

\item M. Kauers and P. Paule, \emph{The Concrete Tetrahedron}, Springer, 2011; C-finite sequences and their closure properties.

\item J. Berstel and C. Reutenauer, \emph{Noncommutative Rational Series with Applications}, Cambridge University Press, 2011; linear representations of counting automata and their reduction.

\end{enumerate}
\end{document}
